% ================================================================
\section{Programação dinâmica}\label{sec:pd}
% ================================================================
\textbf{Motivação (slides).} Na decomposição recursiva, \textbf{subproblemas idênticos} podem aparecer muitas vezes; resolvendo-os de novo a cada aparição, o algoritmo pode ser \textbf{exponencial} mesmo que o número de subproblemas \emph{distintos} seja polinomial. \textbf{PD}: uma tabela com uma entrada por subproblema distinto; cada um é resolvido uma única vez, e as próximas consultas só leem a tabela.
\begin{itemize}
  \item \textbf{D\&C $\times$ PD:} D\&C é bom quando os subproblemas são \textbf{disjuntos} (mergesort); quando eles \textbf{compartilham} subproblemas, troca-se o \emph{top-down} por \emph{bottom-up}: resolver dos menores para os maiores, de modo que, ao visitar um subproblema, todos de que ele depende já estejam prontos.
  \item \textbf{Subestrutura ótima:} uma solução ótima é composta de soluções ótimas de subproblemas (prova típica: ``recorta e cola'' --- se a parte não fosse ótima, troque-a por uma melhor e obtenha uma solução global melhor, contradição).
  \item \textbf{Sobreposição de subproblemas:} a recursão ingênua revisita os mesmos subproblemas (é isso que faz a tabela valer a pena).
  \item \textbf{Memorização} (\emph{top-down}): a recursão original + ``se já calculei, devolvo da tabela''. \textbf{Bottom-up}: laços na ordem certa. Mesma complexidade; o bottom-up evita pilha de recursão e às vezes permite economizar memória.
\end{itemize}

\begin{pseudo}{Fibonacci: recursão ingênua $\Theta(\varphi^n)$, memorização $O(n)$, bottom-up $O(n)$ e memória $O(1)$}
\begin{algorithmic}
\Function{FibMemo}{$n$} \Comment{$F[\,]$ inicializado com $-1$}
  \If{$F[n]=-1$}
    \If{$n\le1$} \State $F[n]:=n$ \Else \State $F[n]:=$ \Call{FibMemo}{$n-1$} $+$ \Call{FibMemo}{$n-2$} \EndIf
  \EndIf
  \State \Return $F[n]$
\EndFunction
\Function{FibBottomUp}{$n$}
  \State $a:=0$;\ $b:=1$ \Comment{$a=F_{i}$, $b=F_{i+1}$}
  \For{$i:=1$ \textbf{até} $n$} \State $(a,b):=(b,\ a+b)$ \EndFor
  \State \Return $a$
\EndFunction
\end{algorithmic}
\end{pseudo}

\begin{prova}[Receita para ``resolva por programação dinâmica'' (vale para toda a parte de PD)]
\begin{enumerate}
  \item \textbf{Subproblema em palavras}: ``$C(i,j)=$ tamanho da LCS de $X_i$ e $Y_j$''. Sem essa frase a recorrência não tem sentido.
  \item \textbf{Decomposição}: o que a solução ótima faz com o \emph{último} elemento (ou com a \emph{primeira} decisão)? Cada caso vira um termo da recorrência.
  \item \textbf{Recorrência + casos base}.
  \item \textbf{Complexidade} $=$ (número de subproblemas) $\times$ (custo de cada um): LCS $nm\times O(1)$; cadeia de matrizes $n^2\times O(n)$; mochila $nW\times O(1)$.
  \item \textbf{Ordem de preenchimento}: cada célula só depende de células já calculadas (linha a linha; por diagonais; por $k$ crescente\dots).
  \item \textbf{Onde está a resposta} ($C[n,m]$, $m[1,n]$, $M[W,n]$, $\max_i L[i]$\dots).
  \item \textbf{Reconstrução}: guarde a escolha feita em cada célula ($s[i,j]$, ``seta'') ou recompute-a voltando pela tabela.
\end{enumerate}
\end{prova}

% ------------------------------------------------------------------
\subsection{Subsequência comum máxima -- LCS (aula 14, 29/09; Cormen §15.4)}\label{sec:lcs}
$Z$ é \textbf{subsequência} de $X=x_1\dots x_n$ se seus símbolos aparecem em $X$ na mesma ordem, \textbf{não necessariamente contíguos} (ex.: $6,10,12$ é subsequência de $3,6,9,10,12,14,6$, mas não de $3,9,10,12,14,6$). Dadas $X$ ($|X|=n$) e $Y$ ($|Y|=m$), queremos uma subsequência comum de tamanho máximo. Força bruta: $2^n$ subsequências de $X$.

\textbf{Decomposição} (olhando os últimos símbolos; $X_i=x_1\dots x_i$): se $x_n=y_m$, esse símbolo termina alguma LCS e o resto é uma LCS de $X_{n-1},Y_{m-1}$; se $x_n\ne y_m$, a LCS não usa $x_n$ ou não usa $y_m$. Com $C(i,j)=$ tamanho da LCS de $X_i$ e $Y_j$:
\[
C(i,j)=\begin{cases}0 & i=0\text{ ou }j=0\\ 1+C(i-1,j-1) & x_i=y_j\\ \max\{C(i-1,j),\,C(i,j-1)\} & x_i\ne y_j\end{cases}\qquad\text{resposta: }C(n,m).
\]
A recursão direta é exponencial; há só $(n+1)(m+1)$ subproblemas distintos $\Rightarrow$ PD.

\begin{pseudo}{LCS bottom-up (slides) + reconstrução (monitoria) -- $\Theta(nm)$ tempo e memória}
\begin{algorithmic}
\Function{LCS}{$X,Y$}
  \For{$j:=0$ \textbf{até} $m$} \State $C[0,j]:=0$ \EndFor
  \For{$i:=1$ \textbf{até} $n$} \State $C[i,0]:=0$ \EndFor
  \For{$i:=1$ \textbf{até} $n$}
    \For{$j:=1$ \textbf{até} $m$}
      \If{$x_i=y_j$} \State $C[i,j]:=1+C[i-1,j-1]$
      \Else \State $C[i,j]:=\max\{C[i-1,j],\,C[i,j-1]\}$
      \EndIf
    \EndFor
  \EndFor
  \State $i:=n$;\ $j:=m$;\ $Z:=$ sequência vazia \Comment{reconstrução: $O(n+m)$}
  \While{$i>0$ \textbf{e} $j>0$}
    \If{$x_i=y_j$} \State coloque $x_i$ no \textbf{início} de $Z$;\ $i:=i-1$;\ $j:=j-1$
    \ElsIf{$C[i-1,j]\ge C[i,j-1]$} \State $i:=i-1$
    \Else \State $j:=j-1$
    \EndIf
  \EndWhile
  \State \Return $Z$ \Comment{$|Z|=C[n,m]$}
\EndFunction
\end{algorithmic}
\end{pseudo}

\begin{exemplo}[Monitoria: $X=\mathrm{ABCBDAB}$, $Y=\mathrm{BDCABA}$ --- conferido]\leavevmode
\begin{center}\small
\begin{tabular}{cc|ccccccc}
 & & $j{=}0$ & 1 & 2 & 3 & 4 & 5 & 6\\
 & & & B & D & C & A & B & A\\ \hline
$i{=}0$ & & 0&0&0&0&0&0&0\\
1 & A & 0&0&0&0&1&1&1\\
2 & B & 0&\textbf{1}&1&1&1&2&2\\
3 & C & 0&1&1&\textbf{2}&2&2&2\\
4 & B & 0&1&1&2&2&\textbf{3}&3\\
5 & D & 0&1&2&2&2&3&3\\
6 & A & 0&1&2&2&3&3&\textbf{4}\\
7 & B & 0&1&2&2&3&4&4\\
\end{tabular}
\end{center}
$C[7,6]=4$. A reconstrução (empate $\Rightarrow$ sobe) passa pelas células em negrito: $(7,6)\uparrow(6,6)$ casa \textbf{A} $\to(5,5)\uparrow(4,5)$ casa \textbf{B} $\to(3,4)\leftarrow(3,3)$ casa \textbf{C} $\to(2,2)\leftarrow(2,1)$ casa \textbf{B}. Como cada símbolo entra no \emph{início}, $Z=\mathrm{BCBA}$. As LCS de tamanho 4 são exatamente \textbf{BCAB, BCBA e BDAB}.
\end{exemplo}
Ex.\ dos slides: $X=5,2,1,3,9,10,14$ e $Y=2,3,1,10$: tamanho 3, soluções $2,3,10$ e $2,1,10$. Ex.\ da lista APA (ex.~6): $X=\mathrm{dacdba}$, $Y=\mathrm{cadcda}$: tamanho \textbf{4} (\textrm{acda} ou \textrm{dcda}).
\begin{itemize}
  \item \textbf{Só o tamanho} com memória $O(\min\{n,m\})$: cada linha só usa a anterior (duas linhas). Mas aí não dá para reconstruir voltando na tabela.
  \item \textbf{Todas as LCS} (lista l6, ex.~3): na reconstrução, quando $x_i\ne y_j$ e $C[i-1,j]=C[i,j-1]$, siga \emph{os dois} caminhos (recursão) e elimine repetições; pode haver exponencialmente muitas.
\end{itemize}

% ------------------------------------------------------------------
\subsection{Multiplicação de cadeia de matrizes (aula 15, 02/10; Cormen §15.2)}\label{sec:cadeia}
$M=M_1\times M_2\times\dots\times M_n$, com $M_i$ de dimensões $r_{i-1}\times r_i$ (consecutivas compatíveis). Multiplicar $(p\times q)\cdot(q\times s)$ custa $p\,q\,s$ multiplicações escalares. O produto é \textbf{associativo}: a ordem não muda o resultado, mas muda (muito) o custo. \textbf{Objetivo:} a parentização de custo mínimo.

\begin{exemplo}[Aula de 02/10 --- conferido]\label{ex:cadeia}
$M_1:20\times50$, $M_2:50\times10$, $M_3:10\times1$, $M_4:1\times100$ ($r=(20,50,10,1,100)$).
$M_1(M_2(M_3M_4))$: $1000+50\,000+100\,000=\mathbf{151\,000}$;\quad $((M_1M_2)M_3)M_4$: $10\,000+200+2000=\mathbf{12\,200}$.
\emph{Nenhuma das duas é a ótima}: a PD abaixo dá $(M_1(M_2M_3))M_4$ com $500+1000+2000=\mathbf{3500}$.
\end{exemplo}
\textbf{Força bruta é exponencial}: o número de parentizações de $n$ matrizes é o número de Catalan $C_{n-1}$ ($1,1,2,5,14,42,132,\dots$), que é $\Omega(4^n/n^{3/2})$.

\textbf{Decomposição:} a \emph{última} multiplicação de $M_i\cdots M_j$ é $(M_i\cdots M_k)\times(M_{k+1}\cdots M_j)$ para algum $i\le k<j$; numa solução ótima, as duas metades também estão parentizadas de forma ótima (recorta e cola). Com $m[i,j]=$ número mínimo de multiplicações para calcular $M_i\cdots M_j$:
\[
m[i,j]=\begin{cases}0 & i=j\\[2pt] \min\limits_{i\le k<j}\big\{\,m[i,k]+m[k+1,j]+r_{i-1}\,r_k\,r_j\,\big\} & i<j\end{cases}\qquad\text{resposta: } m[1,n].
\]
$O(n^2)$ subproblemas, cada um testa $O(n)$ valores de $k$ $\Rightarrow$ \textbf{$\Theta(n^3)$ tempo, $\Theta(n^2)$ memória}. Ordem: $m[i,j]$ depende de intervalos \emph{mais curtos} $\Rightarrow$ preencher \textbf{por diagonais} $d=j-i=1,2,\dots,n-1$ (a diagonal principal é zero).

\begin{pseudo}{Cadeia de matrizes (quadro da aula, com a tabela $s$ da lista l6, ex.~2)}
\begin{algorithmic}
\Function{OrdemÓtima}{$r_0,\dots,r_n$}
  \For{$i:=1$ \textbf{até} $n$} \State $m[i,i]:=0$ \EndFor
  \For{$d:=1$ \textbf{até} $n-1$} \Comment{``para cada diagonal''}
    \For{$i:=1$ \textbf{até} $n-d$} \Comment{``para cada linha da diagonal''}
      \State $j:=i+d$;\ $m[i,j]:=\infty$
      \For{$k:=i$ \textbf{até} $j-1$}
        \State $q:=m[i,k]+m[k+1,j]+r_{i-1}\,r_k\,r_j$
        \If{$q<m[i,j]$} \State $m[i,j]:=q$;\ $s[i,j]:=k$ \EndIf
      \EndFor
    \EndFor
  \EndFor
  \State \Return $m[1,n]$ e $s$
\EndFunction
\Procedure{Imprime}{$s,i,j$}
  \If{$i=j$} \State imprima ``$M_i$''
  \Else \State imprima ``(''; \Call{Imprime}{$s,i,s[i,j]$}; \Call{Imprime}{$s,s[i,j]+1,j$}; imprima ``)''
  \EndIf
\EndProcedure
\Function{Multiplica}{$A,s,i,j$} \Comment{faz $A_i\cdots A_j$ na ordem ótima}
  \If{$i=j$} \State \Return $A_i$ \EndIf
  \State \Return \Call{Multiplica}{$A,s,i,s[i,j]$} $\times$ \Call{Multiplica}{$A,s,s[i,j]+1,j$}
\EndFunction
\end{algorithmic}
\end{pseudo}

\begin{center}\small
\begin{tabular}{c|cccc}
$m[i,j]$ ($s[i,j]$) & $j=1$ & 2 & 3 & 4\\ \hline
$i=1$ & 0 & 10\,000 (1) & 1500 (1) & \textbf{3500 (3)}\\
2 & & 0 & 500 (2) & 5500 (3)\\
3 & & & 0 & 1000 (3)\\
4 & & & & 0
\end{tabular}
\end{center}
Tabela do Exemplo~\ref{ex:cadeia}. Como $s[1,4]=3$, a última multiplicação é $(M_1M_2M_3)\cdot M_4$; como $s[1,3]=1$, a primeira parte é $M_1(M_2M_3)$. Logo a ordem ótima é $((M_1(M_2M_3))M_4)$, com custo 3500.
\begin{itemize}
  \item Monitoria: $r=(10,100,5,50)$: $m[1,2]=5000$, $m[2,3]=25\,000$, $m[1,3]=\min\{75\,000,\ 7500\}=\mathbf{7500}$, ordem $((M_1M_2)M_3)$.
  \item Lista l6, ex.~1 (Cormen 15.2-1), $r=(5,10,3,12,\allowbreak5,50,6)$: custo \textbf{2010}, ordem $((M_1M_2)\allowbreak((M_3M_4)\allowbreak(M_5M_6)))$.
  \item Exemplo do Cormen (Fig.~15.5): $r=(30,35,15,5,\allowbreak10,20,25)$, custo 15\,125, ordem $((M_1(M_2M_3))\allowbreak((M_4M_5)M_6))$.
\end{itemize}

% ------------------------------------------------------------------
\subsection{Mochila 0-1 (monitoria de PD)}\label{sec:mochila}
$n$ itens, item $i$ com peso $p_i$ e valor $v_i$ (inteiros); capacidade $W$. Escolher $S\subseteq\{1,\dots,n\}$ com $\sum_{i\in S}p_i\le W$ maximizando $\sum_{i\in S}v_i$ (cada item entra no máximo uma vez).
Com $M(w,i)=$ maior valor usando só os itens $1..i$ com capacidade $w$ (o item $i$ fica de fora, ou entra se couber):
\[
M(w,i)=\begin{cases}0 & i=0\\ M(w,i-1) & p_i>w\\ \max\{M(w,i-1),\ v_i+M(w-p_i,\,i-1)\} & p_i\le w\end{cases}\qquad\text{resposta: }M(W,n).
\]
\begin{pseudo}{Mochila 0-1 -- $\Theta(nW)$ (pseudo-polinomial)}
\begin{algorithmic}
\Function{Mochila}{$p,v,n,W$}
  \For{$w:=0$ \textbf{até} $W$} \State $M[w,0]:=0$ \EndFor
  \For{$i:=1$ \textbf{até} $n$}
    \For{$w:=0$ \textbf{até} $W$}
      \State $M[w,i]:=M[w,i-1]$
      \If{$p_i\le w$ \textbf{e} $v_i+M[w-p_i,i-1]>M[w,i]$} \State $M[w,i]:=v_i+M[w-p_i,i-1]$ \EndIf
    \EndFor
  \EndFor
  \State $w:=W$ \Comment{reconstrução}
  \For{$i:=n$ \textbf{decrescendo até} $1$}
    \If{$M[w,i]\ne M[w,i-1]$} \State o item $i$ está na solução;\ $w:=w-p_i$ \EndIf
  \EndFor
  \State \Return $M[W,n]$
\EndFunction
\end{algorithmic}
\end{pseudo}
\begin{exemplo}[Monitoria: $W=4$; itens $(p,v)=(1,2),(3,7),(2,1)$ --- conferido]\leavevmode
\begin{center}\small
\begin{tabular}{c|cccc}
$w\backslash i$ & 0 & 1 & 2 & 3\\ \hline
0 & 0 & 0 & 0 & 0\\
1 & 0 & 2 & 2 & 2\\
2 & 0 & 2 & 2 & 2\\
3 & 0 & 2 & 7 & 7\\
4 & 0 & 2 & \textbf{9} & \textbf{9}
\end{tabular}
\end{center}
Ótimo: itens 1 e 2 (peso 4, valor \textbf{9}). \emph{O slide da monitoria mostra 7 na última linha; está errado}: $M(4,2)=\max\{M(4,1),\,7+M(1,1)\}=\max\{2,\,9\}=9$.
\end{exemplo}
\begin{itemize}
  \item Monitoria: $W=50$, $(p,v)=(10,60),(20,100),(30,120)$: ótimo $\mathbf{220}$ (itens 2 e 3). O guloso pela razão $v/p$ ($6,5,4$) pega os itens 1 e 2 e obtém só 160 --- guloso \textbf{não} resolve a mochila 0-1 (resolve a mochila \emph{fracionária}: 240).
  \item Lista APA, ex.~7: $W=100$, $(p,v)=(20,50),(30,70),(40,90),(50,120),(60,140)$: ótimo \textbf{240} com os itens 1, 2 e 4 (peso exatamente 100).
  \item \textbf{Por que pseudo-polinomial?} O tempo $\Theta(nW)$ depende do \emph{valor} de $W$, mas a entrada guarda $W$ em $\log_2W$ bits: com $W=10^9$ (30 bits) o algoritmo faz $10^9$ iterações por item. É exponencial no tamanho da entrada (a mochila é NP-difícil).
  \item \textbf{Memória $O(W)$} (MC458, ex.~4): um único vetor $K[0..W]$; para cada item, \textbf{$w$ decrescendo} de $W$ até $p_i$: $K[w]:=\max\{K[w],K[w-p_i]+v_i\}$. Com $w$ \emph{crescendo}, o item poderia ser reutilizado: vira a mochila \emph{ilimitada}.
\end{itemize}

% ------------------------------------------------------------------
\subsection{Lista de PD resolvida}\label{sec:listapd}
\begin{enumerate}
  \item \textbf{Cadeia de matrizes \emph{maximizando} o custo tem subestrutura ótima?} \textbf{Sim.} Se a última multiplicação de uma parentização ótima (máxima) de $M_i\cdots M_j$ é no ponto $k$, o custo é $c(i..k)+c(k+1..j)+r_{i-1}r_kr_j$, e o último termo depende só de $k$, não de como cada metade foi parentizada. Se uma metade não fosse máxima, trocá-la por uma máxima aumentaria o total: contradição. Recorrência igual, com $\max$; $\Theta(n^3)$.
  \item \textbf{Maior subsequência crescente em $O(n^2)$.} $L[i]=$ tamanho da maior subsequência crescente \emph{terminando em} $a_i$: $L[i]=1+\max\{L[j]: j<i,\ a_j<a_i\}$ ($\max\emptyset=0$); resposta $\max_iL[i]$; guarde $ant[i]$ (o $j$ escolhido) para reconstruir. Ex.: $5,2,8,6,3,6,9,7$: $L=1,1,2,2,2,3,4,4$; resposta 4 ($2,3,6,9$). \emph{Alternativa:} LCS entre $a$ e $a$ ordenado sem repetições. \emph{$O(n\log n)$} (l6, ex.~4): mantenha $d[\ell]=$ menor final possível de uma subsequência crescente de tamanho $\ell$ (vetor crescente); para cada $a_i$, busca binária pelo primeiro $d[\ell]\ge a_i$ e substitua (ou acrescente no fim).
  \item \textbf{Maior subsequência palíndroma.} $P(i,j)$ para o trecho $s_i..s_j$: $P(i,i)=1$, $P(i,i-1)=0$; se $s_i=s_j$, $P(i,j)=2+P(i+1,j-1)$; senão $\max\{P(i+1,j),P(i,j-1)\}$. Preenche por tamanho do trecho; $\Theta(n^2)$. ``character'' $\to$ 5 (\textrm{carac}). (Equivale ao tamanho da LCS entre $s$ e $s$ invertida.)
  \item \textbf{Caminho mais longo de $s$ a $t$ num DAG com pesos.} Ordene topologicamente (§\ref{sec:topo}); $D(s)=0$, demais $-\infty$; na ordem topológica, $D(v)=\max_{uv\in E}\{D(u)+w(u,v)\}$; guarde o predecessor. $\Theta(n+m)$. Num DAG todo caminho é simples, por isso a subestrutura funciona; em grafos gerais, caminho simples mais longo é NP-difícil (contém o caminho hamiltoniano).
  \item \textbf{Mochila inteira (ilimitada).} $K(w)=\max\big\{0,\ \max_{i:\,p_i\le w}\{K(w-p_i)+v_i\}\big\}$, $w=0..W$; $\Theta(nW)$. Ex.: $(p,v)=(3,5),(4,7),(2,3)$, $W=10$: 17 (p.ex.\ $4+4+2$).
  \item \textbf{Cortar uma string (mesma ideia da cadeia de matrizes).} Ordene os pontos e acrescente sentinelas $L_0=0$ e $L_{m+1}=n$. $C(i,j)=$ custo mínimo de fazer todos os cortes estritamente entre $L_i$ e $L_j$: $C(i,i+1)=0$ e
  \[C(i,j)=(L_j-L_i)+\min_{i<k<j}\{C(i,k)+C(k,j)\}\quad(\text{o primeiro corte do trecho é em }L_k).\]
  Resposta $C(0,m+1)$; guarde o $k$ ótimo para listar a sequência de cortes (pré-ordem). $\Theta(m^3)$. Ex.\ do enunciado ($n=20$, cortes $2,8,10$): ótimo \textbf{38}, cortando primeiro em 10 e depois em 2 e 8 (qualquer ordem); as seis ordens custam $50,46,40,40,38,38$. (A ``tora'' da l6, ex.~14, é o mesmo problema: $10$ m com cortes $2,4,7$ $\Rightarrow$ 20.)
\end{enumerate}

\begin{cuidado}[Armadilhas de PD]
\begin{itemize}
  \item \textbf{Subsequência $\ne$ subvetor}: LCS não exige contiguidade; o ``subvetor de soma máxima'' (§\ref{sec:dc}) exige.
  \item Recorrência sem caso base, ou sem dizer \emph{o que} $C(i,j)$ significa.
  \item Ordem de preenchimento errada: na cadeia de matrizes, preencher por \emph{linhas} $i=1,2,\dots$ usa $m[i+1,j]$ ainda não calculado. Vá por diagonais (ou $i$ decrescente).
  \item Cadeia de matrizes: o termo é $r_{i-1}\,r_k\,r_j$ (não $r_i r_k r_j$); e $k$ vai de $i$ a $j-1$.
  \item Dizer que $O(nW)$ é polinomial. Não é: é pseudo-polinomial.
  \item Material da monitoria: a tabela da mochila ($M(4,2)=9$, não 7) e a lista de LCS de ABCBDAB/BDCABA (são três: BCAB, \textbf{BCBA}, BDAB --- e a reconstrução mostrada devolve justamente BCBA, já na ordem certa, porque insere no início).
\end{itemize}
\end{cuidado}

\begin{chave}
PD $=$ recorrência com subproblemas que se repetem $+$ tabela. LCS $\Theta(nm)$; cadeia de matrizes $\Theta(n^3)$ por diagonais; mochila 0-1 $\Theta(nW)$ pseudo-polinomial; Floyd-Warshall $\Theta(n^3)$ (§\ref{sec:fw}). Sempre: significado do subproblema, recorrência, base, ordem, complexidade, reconstrução.
\end{chave}
